% !TEX TS-program = LuaLaTeX

% sansmathfonts.tex
%
% Documentation for the sansmathfonts package
%
% author: Ariel Barton
%
% Copyright Ariel Barton, 2013--2024
%
% This work may be distributed and/or modified under the
% conditions of the LaTeX Project Public License, either
% version 1.3c of this license or (at your option) any
% later version.
% The latest version of the license is in
%    http://www.latex-project.org/lppl.txt
% and version 1.3 or later is part of all distributions of
% LaTeX version 2003/06/01 or later.
%
% This work has the LPPL maintenance status "author-maintained".
%
% The complete list of files considered part of this work is in
% the file `sansmathfonts.pdf' and its source code `sansmathfonts.tex'.
%
% Date: 2024/10/15

\documentclass{amsart}

\raggedbottom

\usepackage[OT1]{fontenc}
\usepackage[I,AMS,letters]{sansmathfonts}
%\renewcommand{\sfdefault}{xcmss}\let\mathserif\mathrm
\usepackage{hyperref, amsfonts, amssymb, esint, slantsc, bm, enumitem}

%\let\mathserif\mathrm

\typeout{********************************}
\typeout{Current font encoding: \csname f@encoding\endcsname}

\def\sectionautorefname{Section}

\title{The \textsf{sansmathfonts} package}
\author{Ariel Barton}

\DeclareTextFontCommand{\cmsmf}
{\fontencoding{OT1}\fontfamily{cmsmf}\selectfont}

\begin{document}



\maketitle



The Computer Modern font family has a sans serif typeface. However, compared to the serif typeface, it is incomplete: there are no sans serif small caps or math fonts. Furthermore, the bold slanted font is not available as an outline font.
This leads to highly unsatisfactory typography of documents that use sans serif for the body text.


The \textsf{sansmathfonts} package provides these ``missing'' fonts. Most of the usefulness of the package is in the fonts; \texttt{sansmathfonts.sty} is a small package providing \LaTeX\ support. To use it, simply say \texttt{\char`\\usepackage\char`\{sansmathfonts\char`\}} in the document preamble.

In the default (OT1) text font encoding, and also in the T1 and U font encodings, this will redefine the document's default sans serif font family from \textsf{cmss} to \textsf{xcmss}; this will make the \textsf{\textbf{\textit{bold slanted}}} and \textsf{\textsc{Caps and small caps}} fonts available via normal \LaTeX\ font commands (\verb|\textbf|, \verb|\textit| and \verb|\textsc|). If you additionally load Harald Harder's \textsf{slantsc} package, this will make \textsf{\textit{\textsc{slanted caps and small caps}}} available.

This will also switch the math fonts to sans serif:
\[\Im \mathop{\mathrm{exp}}(i\omega)=\sin(\omega) \]
If you use symbols from the \textsf{amsfonts} or \textsf{esint} packages, they will also be replaced by appropriate sans serif versions:
\[\ointclockwise \quad \mho \quad \hbar\]
By default, the commands \verb|\mathrm| and \verb|\mathsf| both produce sans serif text. To get serifed roman text, use the command \verb|\mathserif|:
\[\mathrm{mathrm}\quad\mathsf{mathsf}\quad\mathserif{mathserif}\]
\textsf{sansmathfonts} knows about the \textsf{beamer} document class and will automatically use \textsf{beamer}'s \texttt{professionalfonts} theme.

The math fonts differ slightly from Knuth's standard sans serif fonts. Specifically, for ease of reading I have chosen to put (some of) the serifs back on the uppercase I, Pi and~Xi:
\[{I}\quad\mathrm{I}\quad\mathbf{I}\quad\bm{I}\quad
\mathit{\Pi}\quad{\Pi}\quad\mathbf{\Pi}\quad
\mathit{\Xi}\quad{\Xi}\quad\mathbf{\Xi}\quad
\text{and not\quad \textsf{\font \rawfont = cmss10 \rawfont I\quad \char5\quad\char4}}\]
\textsf{Sans serif Is outside of math mode} still have no serifs unless the package option \texttt{[I]} is used; note that this option as yet only works with the OT1 and U font encodings.


Feedback is appreciated and may be sent to \texttt{origamist@gmail.com}.

\section{Package options}

\begin{itemize}[leftmargin=*]
\item \texttt{[I]} The \texttt{[I]} package option puts the serifs back on the capital I even in text mode. This option only works with the OT1 and U font encodings. It will work under \textsf{pdflatex}'s defaults; in Lua\LaTeX\ or Xe\LaTeX, you will need to change the text encoding by saying \texttt{\string\usepackage[OT1]\char`\{fontenc\char`\}}.
\item \texttt{[onlymath]}, \texttt{[nottext]}. These options provide sans serif math but do not change the text sans serif font.
\item \texttt{[onlytext]}, \texttt{[notmath]}. These options provide sans serif text fonts and improve the behavior of \texttt{\char`\\mathsf} but do not change the default math font from roman to sans serif. You can get a similar effect by not using the \textsf{sansmathfonts} package and instead using the lines
\\\null\qquad\texttt{\char`\\renewcommand\char`\{\char`\\sfdefault\char`\}\char `\{xcmss\char`\}}
\\\null\qquad\texttt{\char`\\DeclareFontFamilySubstitution\char`\{TS1\char`\}\char`\{xcmss\char`\}\char`\{cmss\char`\}}
\\
or
\\\null\qquad
\texttt{\char`\\renewcommand\char`\{\char`\\sfdefault\char`\}\char `\{cmsmf\char`\}}
\\\null\qquad\texttt{\char`\\DeclareFontFamilySubstitution\char`\{TS1\char`\}\char`\{cmsmf\char`\}\char`\{cmss\char`\}}
\\
in the document preamble.
\item \texttt{[AMS]} This package option defines the following five commands:

\begin{tabular}{lc}
\texttt{\char`\\sanshbar}&$\sanshbar$\\
\texttt{\char`\\sanshslash}&$\sanshslash$\\
\texttt{\char`\\sansmho}&$\sansmho$\\
\texttt{\char`\\sanseth}&$\sanseth$\\
\texttt{\char`\\sansbackepsilon}&$\sansbackepsilon$\\
\end{tabular}

If the \texttt{[onlytext]} and \texttt{[notmath]} options are not selected, these commands are simply aliases for the corresponding AMS symbols \texttt{\char`\\hbar}, \texttt{\char`\\hslash}, \texttt{\char`\\mho}, \texttt{\char`\\eth}, \texttt{\char`\\backepsilon} (and will only work if the \textsf{amssymb} package is loaded). 

However, if the \texttt{[onlytext]} or \texttt{[notmath]} option is selected, then the above five commands will be redefined to produce sans serif symbols, even though \texttt{\char`\\hbar} and so on produce serifed symbols.

\item\texttt{[letters]} This package option is similar to the \texttt{[AMS]} option. It defines sans serif versions of all the math symbols included in \TeX's letters font:

\begin{tabular}[t]{lc}
\texttt{\char`\\sansalpha}&$\sansalpha$\\
\texttt{\char`\\sansbeta}&$\sansbeta$\\
\texttt{\char`\\sansgamma}&$\sansgamma$\\
\texttt{\char`\\sansdelta}&$\sansdelta$\\
\texttt{\char`\\sansepsilon}&$\sansepsilon$\\
\texttt{\char`\\sanszeta}&$\sanszeta$\\
\texttt{\char`\\sanseta}&$\sanseta$\\
\texttt{\char`\\sanstheta}&$\sanstheta$\\
\texttt{\char`\\sansiota}&$\sansiota$\\
\texttt{\char`\\sanskappa}&$\sanskappa$\\
\texttt{\char`\\sanslambda}&$\sanslambda$\\
\texttt{\char`\\sansmu}&$\sansmu$\\
\texttt{\char`\\sansnu}&$\sansnu$\\
\end{tabular}
\begin{tabular}[t]{lc}
\texttt{\char`\\sansxi}&$\sansxi$\\
\texttt{\char`\\sanspi}&$\sanspi$\\
\texttt{\char`\\sansrho}&$\sansrho$\\
\texttt{\char`\\sanssigma}&$\sanssigma$\\
\texttt{\char`\\sanstau}&$\sanstau$\\
\texttt{\char`\\sansupsilon}&$\sansupsilon$\\
\texttt{\char`\\sansphi}&$\sansphi$\\
\texttt{\char`\\sanschi}&$\sanschi$\\
\texttt{\char`\\sanspsi}&$\sanspsi$\\
\texttt{\char`\\sansomega}&$\sansomega$\\
\texttt{\char`\\sansvarepsilon}&$\sansvarepsilon$\\
\texttt{\char`\\sansvartheta}&$\sansvartheta$\\
\texttt{\char`\\sansvarpi}&$\sansvarpi$\\
\end{tabular}
\begin{tabular}[t]{lc}
\texttt{\char`\\sansvarrho}&$\sansvarrho$\\
\texttt{\char`\\sansvarsigma}&$\sansvarsigma$\\
\texttt{\char`\\sansvarphi}&$\sansvarphi$\\
\texttt{\char`\\sansstar}&$\sansstar$\\
\texttt{\char`\\sanspartial}&$\sanspartial$\\
\texttt{\char`\\sansflat}&$\sansflat$\\
\texttt{\char`\\sansnatural}&$\sansnatural$\\
\texttt{\char`\\sanssharp}&$\sanssharp$\\
\texttt{\char`\\sanssmile}&$\sanssmile$\\
\texttt{\char`\\sansfrown}&$\sansfrown$\\
\texttt{\char`\\sansell}&$\sansell$\\
\texttt{\char`\\sanswp}&$\sanswp$\\
\end{tabular}

As with \texttt{[AMS]}, these are aliases to the corresponding standard commands if \texttt{[onlytext]} and \texttt{[notmath]} are not selected, and are a new math alphabet if \texttt{[onlytext]} or \texttt{[notmath]} is selected.

Note that sans symbols for uppercase Greek letters are \emph{not} provided, as \texttt{\char`\\mathsf\char`\{\char`\\Gamma\char`\}} produces $\mathsf{\Gamma}$ (not \char0) even if \texttt{[onlytext]} or \texttt{[notmath]} is selected.


\end{itemize}



\section{List of new fonts}\label{sec:fonts}

All of the Type 1 fonts in this package were generated using mftrace 1.2.18 and Fontforge.

The following fonts are based mainly on Donald Knuth's Computer Modern fonts.

\textsf{\fontshape{ui}\selectfont Unslanted italic} (needed in some versions of \TeX\ for the pounds symbol \textsf{\pounds}):

\nobreak

\noindent\hskip0.06\textwidth\parbox[t]{0.235\textwidth}{
\begin{itemize}[leftmargin=*]
\item \textsf{cmssu10}
\end{itemize}}

\bigskip

Text \textsf{\textsc{caps and small caps}}, OT1 encoding:

\nobreak

\noindent\hskip0.06\textwidth\parbox[t]{0.235\textwidth}{
\begin{itemize}[leftmargin=*]
\item \textsf{cmssbxcsc10}
\item \textsf{cmssxicsc10}
\end{itemize}}\parbox[t]{0.235\textwidth}{\begin{itemize}[leftmargin=*]
\item \textsf{cmsscsc8}
\item \textsf{cmsscsci8}
\end{itemize}}\parbox[t]{0.235\textwidth}{\begin{itemize}[leftmargin=*]
\item \textsf{cmsscsc9}
\item \textsf{cmsscsci9}
\end{itemize}}\parbox[t]{0.235\textwidth}{\begin{itemize}[leftmargin=*]
\item \textsf{cmsscsc10}
\item \textsf{cmsscsci10}
\end{itemize}}

\bigskip

Math italic ($\alpha \beta \gamma abc \ell \wp$):

\nobreak

\noindent\hskip0.06\textwidth\parbox[t]{0.235\textwidth}{
\begin{itemize}[leftmargin=*]
\item \textsf{cmssmi5}
\item \textsf{cmssmi6}
\item \textsf{cmssmi7}
\end{itemize}}\parbox[t]{0.235\textwidth}{\begin{itemize}[leftmargin=*]
\item \textsf{cmssmi8}
\item \textsf{cmssmi9}
\item \textsf{cmssmi10}
\end{itemize}}\parbox[t]{0.235\textwidth}{\begin{itemize}[leftmargin=*]
\item \textsf{cmssmib5}
\item \textsf{cmssmib6}
\item \textsf{cmssmib7}
\end{itemize}}\parbox[t]{0.235\textwidth}{\begin{itemize}[leftmargin=*]
\item \textsf{cmssmib8}
\item \textsf{cmssmib9}
\item \textsf{cmssmib10}
\end{itemize}}

\bigskip

Math symbols ($\Re\oplus \Im  $):

\nobreak

\noindent\hskip0.06\textwidth\parbox[t]{0.235\textwidth}{
\begin{itemize}[leftmargin=*]
\item \textsf{cmsssy5}
\item \textsf{cmsssy6}
\item \textsf{cmsssy7}
\end{itemize}}\parbox[t]{0.235\textwidth}{\begin{itemize}[leftmargin=*]
\item \textsf{cmsssy8}
\item \textsf{cmsssy9}
\item \textsf{cmsssy10}
\end{itemize}}\parbox[t]{0.235\textwidth}{\begin{itemize}[leftmargin=*]
\item \textsf{cmssbsy5}
\item \textsf{cmssbsy6}
\item \textsf{cmssbsy7}
\end{itemize}}\parbox[t]{0.235\textwidth}{\begin{itemize}[leftmargin=*]
\item \textsf{cmssbsy8}
\item \textsf{cmssbsy9}
\item \textsf{cmssbsy10}
\end{itemize}}

\bigskip


Math extended fonts ($\int \sum \prod$):

\nobreak

\noindent\hskip0.06\textwidth\parbox[t]{0.235\textwidth}{
\begin{itemize}[leftmargin=*]
\item \textsf{cmssex7}
\end{itemize}}\parbox[t]{0.235\textwidth}{\begin{itemize}[leftmargin=*]
\item \textsf{cmssex8}
\end{itemize}}\parbox[t]{0.235\textwidth}{\begin{itemize}[leftmargin=*]
\item \textsf{cmssex9}
\end{itemize}}\parbox[t]{0.235\textwidth}{\begin{itemize}[leftmargin=*]
\item \textsf{cmssex10}
\end{itemize}}

\bigskip

\cmsmf{Sans serif text fonts with serifed capital~I:}

\noindent\hskip0.06\textwidth\parbox[t]{0.235\textwidth}{
\begin{itemize}[leftmargin=*]
\item \textsf{cmsmf8}
\item \textsf{cmsmf9}
\item \textsf{cmsmf10}
\item \textsf{cmsmf12}
\item \textsf{cmsmf17}
\item \textsf{cmsmfcsc8}
\item \textsf{cmsmfcsc9}
\item \textsf{cmsmfcsc10}
\end{itemize}\par}\parbox[t]{0.235\textwidth}
{\begin{itemize}[leftmargin=*]
\item \textsf{cmsmfbx8}
\item \textsf{cmsmfbx9}
\item \textsf{cmsmfbx10}
\item \textsf{cmsmfbx12}
\item \textsf{cmsmfbx17}
\item \textsf{cmsmfbxcsc10}
\end{itemize}}\parbox[t]{0.235\textwidth}
{\begin{itemize}[leftmargin=*]
\item \textsf{cmsmfi8}
\item \textsf{cmsmfi9}
\item \textsf{cmsmfi10}
\item \textsf{cmsmfi12}
\item \textsf{cmsmfi17}
\item \textsf{cmsmfcsci8}
\item \textsf{cmsmfcsci9}
\item \textsf{cmsmfcsci10}
\end{itemize}}\parbox[t]{0.235\textwidth}
{\begin{itemize}[leftmargin=*]
\item \textsf{cmsmfxi8}
\item \textsf{cmsmfxi9}
\item \textsf{cmsmfxi10}
\item \textsf{cmsmfxi12}
\item \textsf{cmsmfxi17}
\item \textsf{cmsmfxicsc10}
\end{itemize}}


\bigskip


The following fonts are based on fonts by other authors.

\nobreak

\noindent\hskip0.06\textwidth\parbox[t]{0.31\textwidth}{
\raggedright
Eddie Saudrais's \textsf{esint} package
\begin{itemize}[leftmargin=*]
\item \textsf{ssesint7}
\item \textsf{ssesint8}
\item \textsf{ssesint9}
\item \textsf{ssesint10}
\end{itemize}}\parbox[t]{0.31\textwidth}{
\raggedright
AMS symbols (\textsf{amsfonts} package)
\begin{itemize}[leftmargin=*]
\item \textsf{ssmsam5}
\item \textsf{ssmsam6}
\item \textsf{ssmsam7}
\item \textsf{ssmsam8}
\item \textsf{ssmsam9}
\item \textsf{ssmsam10}
\end{itemize}}\parbox[t]{0.31\textwidth}{
\raggedright
AMS symbols (\textsf{amsfonts} package)
\begin{itemize}[leftmargin=*]
\item \textsf{ssmsbm5}
\item \textsf{ssmsbm6}
\item \textsf{ssmsbm7}
\item \textsf{ssmsbm8}
\item \textsf{ssmsbm9}
\item \textsf{ssmsbm10}
\end{itemize}}

\bigskip

The following fonts are based on J\"org Knappen's European Computer Modern fonts.

\nobreak

\bigskip
\noindent\hskip0.06\textwidth
\parbox[t]{0.235\textwidth}{
\font\temp = eczz1000 \temp
Normal
\begin{itemize}[leftmargin=*]
\item \textsf{eczz0500}
\item \textsf{eczz0600}
\item \textsf{eczz0700}
\item \textsf{eczz0800}
\item \textsf{eczz0900}
\item \textsf{eczz1000}
\item \textsf{eczz1095}
\item \textsf{eczz1200}
\item \textsf{eczz1440}
\item \textsf{eczz1728}
\item \textsf{eczz2074}
\item \textsf{eczz2488}
\item \textsf{eczz2986}
\item \textsf{eczz3583}
\end{itemize}
}\parbox[t]{0.235\textwidth}{
\font\temp = eczi1000 \temp
Slanted
\begin{itemize}[leftmargin=*]
\item \textsf{eczi0500}
\item \textsf{eczi0600}
\item \textsf{eczi0700}
\item \textsf{eczi0800}
\item \textsf{eczi0900}
\item \textsf{eczi1000}
\item \textsf{eczi1095}
\item \textsf{eczi1200}
\item \textsf{eczi1440}
\item \textsf{eczi1728}
\item \textsf{eczi2074}
\item \textsf{eczi2488}
\item \textsf{eczi2986}
\item \textsf{eczi3583}
\end{itemize}
}\parbox[t]{0.235\textwidth}{
\font\temp = eczx1000 \temp
Bold
\begin{itemize}[leftmargin=*]
\item \textsf{eczx0500}
\item \textsf{eczx0600}
\item \textsf{eczx0700}
\item \textsf{eczx0800}
\item \textsf{eczx0900}
\item \textsf{eczx1000}
\item \textsf{eczx1095}
\item \textsf{eczx1200}
\item \textsf{eczx1440}
\item \textsf{eczx1728}
\item \textsf{eczx2074}
\item \textsf{eczx2488}
\item \textsf{eczx2986}
\item \textsf{eczx3583}
\end{itemize}
}\parbox[t]{0.235\textwidth}{
\font\temp = eczo1000 \temp
Bold slanted
\begin{itemize}[leftmargin=*]
\item \textsf{eczo0500}
\item \textsf{eczo0600}
\item \textsf{eczo0700}
\item \textsf{eczo0800}
\item \textsf{eczo0900}
\item \textsf{eczo1000}
\item \textsf{eczo1095}
\item \textsf{eczo1200}
\item \textsf{eczo1440}
\item \textsf{eczo1728}
\item \textsf{eczo2074}
\item \textsf{eczo2488}
\item \textsf{eczo2986}
\item \textsf{eczo3583}
\end{itemize}
}

\bigskip

The following fonts are part of the \textsf{sauter} package and are supplied with Mac\TeX\ 2012 as Metafont (\texttt{.mf}) fonts. These provide \textsf{\textbf{bold}} and \textsf{\textsl{\textbf{bold slanted}}} fonts at varying sizes.
The \textsf{sansmathfonts} package provides outline versions.

\noindent\hskip0.06\textwidth\parbox[t]{0.235\textwidth}{
\font\temp = eczo1000 \temp
\begin{itemize}[leftmargin=*]
\item \textsf{cmssxi8}
\item \textsf{cmssxi9}
\item \textsf{cmssxi10}
\end{itemize}
}\parbox[t]{0.235\textwidth}{
\begin{itemize}[leftmargin=*]
\item \textsf{cmssxi12}
\item \textsf{cmssxi17}
\end{itemize}
}\parbox[t]{0.235\textwidth}{
\begin{itemize}[leftmargin=*]
\item \textsf{cmssbx8}
\item \textsf{cmssbx9}
\end{itemize}
}\parbox[t]{0.235\textwidth}{
\begin{itemize}[leftmargin=*]
\item \textsf{cmssbx12}
\item \textsf{cmssbx17}
\end{itemize}
}


\section{Files in this package}\label{sec:files}

\subsection{New fonts}
109 of the 146 new fonts listed in \autoref{sec:fonts} come in three files each: the \TeX\ Font Metric files (extension \texttt{.tfm}), the Type 1 font file (extension \texttt{.pfb}), and Metafont source file (extension \texttt{.mf}).

\subsection{Virtual fonts}
The 28 \textsf{cmsmf} fonts are almost identical to their \textsf{cmss} counterparts. Thus, these fonts are provided as \emph{virtual} fonts, and so come in five parts: the virtual font file (\texttt{cmsmf10.vf}), the \TeX\ Font Metric file (\texttt{cmsmf10.tfm}), and the font \textsf{cmsmfIPiXi10} containing only the altered letters \cmsmf{I}, \cmsmf{\char4} and \cmsmf{\char5} (and \cmsmf{\textsc{i}}, in the small caps fonts); this font comes as MetaFont source (\texttt{cmsmfIPiXi10.mf}), \TeX\ font metric (\texttt{cmsmfIPiXi10.tfm}) and Type 1 font (\texttt{cmsmfIPiXi10.pfb}).

\subsection{Outline versions of preexisting fonts}
The Metafont source for the 9 \textsf{cmssxi} and \textsf{cmssbx} fonts are part of the \textsf{sauter} package. As such, I have not included any new \texttt{.mf} files for these fonts. I have provided outline versions (\texttt{.pfb} files) of these 9 fonts, as without the \texttt{.pfb} files these fonts are provided as bitmaps only, which looks very ugly on most modern displays.

I have also not included the files \texttt{cmssxi10.tfm}, etc., because they may be automatically generated from the corresponding \texttt{.mf} files and may already be included in some distributions.

However, in order to allow compilation on systems without the \textsf{sauter} package, I wanted to include \texttt{.tfm} files for these 9 fonts. For the reasons given above, I cannot title these fonts with the expected name of \texttt{cmssxi10.tfm}; instead I have provided a file called \texttt{cmssxi10-cmsmfcopy.tfm} and included instructions (in the \texttt{.map} file) stating that both \texttt{cmssxi10-cmsmfcopy.tfm} and \texttt{cmssxi10.tfm} should use the outline font \texttt{cmssxi10.pfb}.

\subsection{Metafont source}
In addition, this package should come with the following 29 supplementary Metafont source files:

\begin{itemize}[leftmargin=*]
\item \texttt{eczi.mf}
\item \texttt{eczo.mf}
\item \texttt{eczx.mf}
\item \texttt{eczz.mf}
\item \texttt{sans-amsya.mf}
\item \texttt{sans-amsyb.mf}
\item \texttt{sans-asymbols.mf}
\item \texttt{sans-bigdel.mf}
\item \texttt{sans-bigint.mf}
\item \texttt{sans-bigop.mf}
\item \texttt{sans-bsymbols.mf}
\item \texttt{sans-calu.mf}
\item \texttt{sans-csc.mf}
\item \texttt{sans-greekl.mf}
\item \texttt{sans-greeku.mf}
\item \texttt{sans-IPiXi.mf}
\item \texttt{sans-IPiXicsc.mf}
\item \texttt{sans-mathex.mf}
\item \texttt{sans-mathint.mf}
\item \texttt{sans-mathsl.mf}
\item \texttt{sans-mathsy.mf}
\item \texttt{sans-roman.mf}
\item \texttt{sans-romanu.mf}
\item \texttt{sans-romms.mf}
\item \texttt{sans-slantms.mf}
\item \texttt{sans-sym.mf}
\item \texttt{sans-symbol.mf}
\item \texttt{sans-xbbold.mf}
\item \texttt{sansfontbase.mf}
\end{itemize}

This package should also come with the following 11 \LaTeX\ Font Definition files:

\begin{itemize}[leftmargin=*]
\item \texttt{omlcmssm.fd}
\item \texttt{omscmsssy.fd}
\item \texttt{omxcmssex.fd}

\item \texttt{ot1cmsmf.fd}
\item \texttt{ot1xcmss.fd}
\item \texttt{t1xcmss.fd}

\item \texttt{ucmsmf.fd}
\item \texttt{ussesint.fd}
\item \texttt{ussmsa.fd}
\item \texttt{ussmsb.fd}
\item \texttt{uxcmss.fd}
\end{itemize}

Finally, it should come with the font map file, LaTeX package, and documentation:

\begin{itemize}[leftmargin=*]
\item \texttt{sansmathfonts.map}
\item \texttt{sansmathfonts.sty}
\item \texttt{sansmathfonts.tex}
\item \texttt{sansmathfonts.pdf}
\end{itemize}

% The subfolders should contain:

% doc: 2 files
% map: 1 file
% tex: 12 files

% New Type 1 fonts: 	137
% Old fonts as Type 1: 	9
% Virtual fonts:		28
% Supplementary .mf:	29

% source: 166 files (=137+29)
% tfm: 174 files	(=137+28+9)
% type1: 146 files  (=137+9)
% vf: 28  			(=28)

\section{License}

This work (the \textsf{sansmathfonts} package) consists of the
files listed in \autoref{sec:files}.

This work may be distributed and/or modified under the
conditions of the \LaTeX\ Project Public License, either
version 1.3c of this license or (at your option) any
later version.

The latest version of the license is in
\begin{quote}
\href {http://www.latex-project.org/lppl.txt} 
{\texttt{http://www.latex-project.org/lppl.txt}}
\end{quote}
and version 1.3 or later is part of all distributions of
\LaTeX\ version 2003/06/01 or later.

This work has the LPPL maintenance status ``maintained''.

Almost all of the Metafont files in this package are very closely based on existing files in the 2011 \TeX\ Live distribution; see comments near the start of the individual files for notes on their sources. Also, note that the 
files
\begin{itemize}[leftmargin=*]
\item \texttt{cmssxi8.pfb}, \texttt{cmssxi8-cmsmfcopy.tfm}
\item \texttt{cmssxi9.pfb}, \texttt{cmssxi9-cmsmfcopy.tfm}
\item \texttt{cmssxi10.pfb}, \texttt{cmssxi10-cmsmfcopy.tfm}
\item \texttt{cmssxi12.pfb}, \texttt{cmssxi12-cmsmfcopy.tfm}
\item \texttt{cmssxi17.pfb}, \texttt{cmssxi17-cmsmfcopy.tfm}
\item \texttt{cmssbx8.pfb}, \texttt{cmssbx8-cmsmfcopy.tfm}
\item \texttt{cmssbx9.pfb}, \texttt{cmssbx9-cmsmfcopy.tfm}
\item \texttt{cmssbx12.pfb}, \texttt{cmssbx12-cmsmfcopy.tfm}
\item \texttt{cmssbx17.pfb}, \texttt{cmssbx17-cmsmfcopy.tfm}
\end{itemize}
were derived from unedited MetaFont source files in the \textsf{sauter} package using \textsf{mktextfm}, \textsf{mftrace 1.2.18} and Fontforge. 

\section{Revision history}

\begin{itemize}[leftmargin=*]
\item
April 2013: Original upload

\item
February 2017: Corrected the font names in \texttt{sansmathfonts.map}; this allowed the package to be used correctly with \textsf{dvips}.  

\item
April 2019: Fixed a bug in the file \texttt{ucmsmf.fd} that prevented the \texttt{[I]} package option from working correctly; rewrote most of the \texttt{.fd} files to allow fonts to be loaded at arbitrary sizes; changed maintenance status from ``author-maintained'' to ``maintained''.

\item
June 2019: Rewrote the file \texttt{omxcmssex.fd} to allow the math extended characters to be loaded at arbitrary sizes.

\item
June 2021: Rewrote the OT1, T1, and U font definition files to substitute bold-extended fonts for bold fonts as necessary. Added some package errors and warnings if the document font encoding is not supported.

\item
October 2022: Bug fix to allow compatibility with Lua\LaTeX\ and Xe\LaTeX.

\item
November 2023: Corrected the spacing of the lowercase \textsc{i} in the caps and small caps fonts used with the \texttt{[I]} package option.

\item October 2024: 

Added the \texttt{[letters]} and \texttt{[AMS]} options. 

Fixed a bug in the \textsf{cmsmf} fonts; in previous versions, when \LaTeX\ asked for bold non-extended fonts, it would substitute \textsf{xcmss} bold-extended fonts rather than \textsf{cmsmf} bold-extended fonts.

Added two lines in \texttt{sansmathfonts.sty} to allow \LaTeX\ to use \texttt{ts1cmss.fd} for TS1 symbols. This is what in modern \LaTeX\ \emph{actually} allows for a correct text mode \texttt{\char`\\pounds} symbol:  \texttt{\char`\\pounds}~\pounds \space \texttt{\char`\\textsf\char`\{\char`\\pounds\char`\}}~\textsf{\pounds}. (A bug remains in that \texttt{\$\char`\\pounds\$} produces $\pounds$ rather than \textsf{\pounds} in math mode; as a workaround, try \texttt{\char`\$\char`\\text\char`\{\char`\\textsf\char`\{\char`\\pounds\char`\}\char`\}\char`\$}.)

Certain \texttt{.tfm} files are supplied with \TeX\ Live as part of the \textsf{sauter} package; I have added copies of these \texttt{.tfm} files to make the \textsf{cmsmf} font family (the \texttt{[I]} package option) work even if the \textsf{sauter} package is not installed.

\end{itemize}




\end{document}